% Compile-ready section. Requires: enumitem, tabularx, booktabs.

\section{GlyphChain Independent Trust Publication and Clean-Device Verification}
\label{sec:glyphchain-independent-trust-publication}

\noindent\textbf{Status:} The immutable public trust-release builder, offline
verifier pin gate, independent package verifier and bounded clean-device
evidence exercise are implemented and locally verified.  Publication of the
checkpoint package and its expected identifier through genuinely independent
production channels, followed by execution on a separate clean device, remain
production release gates.

\subsection{Purpose and Trust Boundary}

A quorum-finality certificate can prove that sufficient members of a stated
validator set signed an exact GlyphChain ledger commitment.  It cannot prove by
itself that the validator set is the legitimate GlyphChain set.  Similarly, a
checkpoint file can contain its own mathematically correct identifier but cannot
make that identifier trusted merely by declaring it.

Tessaris therefore separates two publication objects:

\begin{enumerate}[leftmargin=*]
  \item the \emph{public trust-release package}, containing the normalized
        validator checkpoint, the offline proof verifier and a digest-bound
        release manifest; and
  \item the \emph{independent checkpoint pin}, consisting of the expected
        64-character checkpoint identifier obtained through a separately
        controlled trusted channel.
\end{enumerate}

The checkpoint's claimed identifier necessarily appears inside the public
package so that it can be identified and indexed.  That copy is informational,
not authoritative.  Trust-anchored verification requires the user or operator
to supply the same identifier independently.  A receipt, peer, checkpoint,
release manifest or verifier package is never allowed to appoint its own trust
anchor.

Examples of acceptable channel separation include a versioned release asset
combined with a pin published in an independently controlled transparency log,
DNSSEC record, signed release announcement or separately administered status
site.  Two paths backed by the same mutable storage account or deployment
credential are not operationally independent.

\subsection{Canonical Validator Checkpoint}

The public checkpoint uses the schema
\texttt{glyphchain.validator-set-checkpoint.v1}.  It contains only public
protocol material:

\begin{itemize}[leftmargin=*]
  \item chain and network identifiers;
  \item checkpoint epoch and effective height;
  \item ordered validator identifiers, voting powers and Ed25519 public keys;
  \item the canonical validator-set root; and
  \item the canonical checkpoint identifier.
\end{itemize}

Validators are normalized and sorted by identifier.  Duplicate or empty
identifiers, non-positive voting power, missing public keys, excessive
membership, invalid network fields and unsupported schemas fail closed.  If
(V) is the canonical ordered validator list, its root is

\[
  r_V = \operatorname{SHA256}(\operatorname{Canonical}(V)).
\]

For checkpoint body (C), excluding its final identifier but including
(r_V), the checkpoint identifier is

\[
  i_C = \operatorname{SHA256}(\operatorname{Canonical}(C)).
\]

Any supplied root or identifier must match the recomputed value.  Downloaded
checkpoint verification additionally requires either an expected checkpoint
identifier or validator-set root obtained out of band; absence of both is an
error rather than an integrity-only success.

\subsection{Immutable Public Trust Release}

The operator command
\texttt{backend/scripts/build\_glyphchain\_trust\_release.py} creates a new,
non-overwriting public release directory.  It refuses symbolic-link inputs and
destinations, invalid release versions, malformed source-manifest digests,
oversized inputs, non-canonical checkpoints and a verifier that lacks the
independent-pin gate.  A pre-existing output directory is never replaced.

The package contains exactly three files:

\begin{center}
\begin{tabularx}{\textwidth}{@{}lX@{}}
\toprule
\textbf{File} & \textbf{Function} \\
\midrule
\texttt{validator-set.glyphchain-checkpoint} & Canonical public validator-set
checkpoint; contains no private validator key. \\
\texttt{tessaris-proof-verifier.html} & Network-free browser verifier for the
document, device signature, signer commitments, transaction inclusion and
quorum certificate. \\
\texttt{TRUST\_RELEASE.json} & Release version, creation time, source release-
manifest digest, checkpoint identity, file lengths, SHA--256 file digests,
privacy statement, channel-separation warning and canonical manifest digest. \\
\bottomrule
\end{tabularx}
\end{center}

The manifest is written last.  An interrupted directory without that manifest
is incomplete and cannot pass verification.  Verification requires an exact
three-file set, rejects symbolic links and unexpected files, bounds every file,
recomputes every digest, re-normalizes the checkpoint and compares it with the
independently supplied pin.  It can additionally bind the package to the exact
source release-manifest SHA--256 digest.

The builder returns the checkpoint identifier as a \emph{candidate} for
independent publication and explicitly marks it as not independently published.
Creating the package and printing the candidate do not satisfy the production
publication gate.

\subsection{Offline Browser Verifier Pin Gate}

The standalone browser verifier performs all operations locally and contains no
network request path.  The user selects the \texttt{.tessaris-proof} receipt,
exact signed text document, optional private evidence certificate and public
checkpoint file.  Trust-anchored finality now requires a fifth input: the
independently obtained checkpoint identifier.

The verifier recomputes the checkpoint's validator-set root and checkpoint
identifier.  Its displayed quorum state is constrained as follows:

\begin{center}
\begin{tabularx}{\textwidth}{@{}lX@{}}
\toprule
\textbf{Displayed state} & \textbf{Verified meaning} \\
\midrule
\texttt{SIGNATURES VALID -- TRUST ANCHOR NOT SUPPLIED} & Validator membership,
votes and quorum are internally consistent, but no checkpoint was supplied. \\
\texttt{CHECKPOINT VALID -- INDEPENDENT PIN NOT SUPPLIED} & The certificate and
checkpoint agree and the checkpoint is internally canonical, but its authority
has not been established independently. \\
\texttt{TRUST-ANCHORED} & The certificate, public checkpoint and independently
entered checkpoint identifier all bind to the same validator set and exact
ledger evidence. \\
\texttt{FAILED} & A document, signature, transaction, Merkle path, certificate,
checkpoint or independent-pin comparison failed. \\
\bottomrule
\end{tabularx}
\end{center}

An absent pin can therefore never produce the trust-anchored state.  A malformed
pin or a mismatch between the independent pin and recomputed checkpoint
identifier produces an explicit failure.

\subsection{Bounded Clean-Device Evidence Exercise}

The command
\texttt{backend/scripts/run\_glyphchain\_clean\_device\_verification.py}
automates the cryptographic exercise using only:

\begin{itemize}[leftmargin=*]
  \item the private \texttt{.tessaris-evidence.zip} supplied by a document party;
  \item the independently downloaded public trust-release directory;
  \item the checkpoint identifier obtained through the second trusted channel;
        and
  \item optionally, the expected source release-manifest digest.
\end{itemize}

The exercise does not contact the originating Tessaris installation, GlyphChain
peer network or either publication channel.  It bounds compressed and expanded
archive sizes and member count, rejects duplicate members, absolute paths,
parent-directory traversal and backslash path ambiguity, and requires exactly
one portable proof receipt plus \texttt{document.txt} and
\texttt{evidence-certificate.json}.  Private signer records are used only in
memory to recompute the commitment root.

The following checks must all pass:

\begin{enumerate}[leftmargin=*]
  \item the document bytes reproduce the committed document digest;
  \item the originating device's Ed25519 signature authenticates the receipt;
  \item the private ordered signer records reproduce the public signer-
        commitment root;
  \item the canonical \texttt{DOCUMENT\_PROOF\_V1} transaction envelope
        reproduces the declared transaction hash and identifier;
  \item the Merkle inclusion path reconstructs the block transaction root;
  \item the version-two quorum certificate binds the same transaction root,
        ledger commitment, block, height, chain and network;
  \item validator signatures meet the weighted two-thirds quorum threshold; and
  \item the certificate validator root matches the checkpoint whose identifier
        matches the independently supplied pin.
\end{enumerate}

The resulting owner-only JSON report is content-free.  It contains the evidence-
package digest, a one-way proof-identifier commitment, trust-release identifiers,
boolean verification results and hashes of any error labels.  It excludes the
document text, signer names, email addresses and private certificate.  The
canonical report carries its own SHA--256 digest, is created exclusively and is
never allowed to overwrite an earlier exercise report.

\subsection{Verified Local Behaviour and Remaining Release Gate}

The focused regression verifies canonical checkpoint export, mandatory
out-of-band pinning, wrong-pin rejection, immutable release construction,
manifest and member tamper detection, unexpected-file rejection, source-
manifest binding, symbolic-link refusal, offline-verifier enforcement,
independent-pin user-interface gating, complete trust-anchored evidence-package
verification, document-substitution detection, archive path-traversal rejection
and tamper-evident content-free exercise reporting.  Nineteen focused checkpoint,
trust-release, standalone-verifier and clean-device tests pass.

These local tests prove that the tooling enforces the intended separation when
given independently sourced inputs.  They do not prove that two production
channels are independently administered or that a separate physical or virtual
device successfully obtained them.  Production remains gated on:

\begin{enumerate}[leftmargin=*]
  \item publishing the immutable checkpoint release through the approved
        versioned release channel;
  \item publishing the expected checkpoint identifier through a genuinely
        independent channel; and
  \item running the implemented exercise on a clean device using those two live
        sources and retaining its digest-bound report with the signed release
        evidence.
\end{enumerate}

Successful trust-anchored verification proves cryptographic consistency of the
supplied evidence against the independently pinned validator set.  It does not
prove legal enforceability, qualified-electronic-signature status, signer
identity beyond the configured policy or uninterrupted public availability.
